% =============================================================================
%  CHAPTER 6 - CONCLUSIONS AND RECOMMENDATIONS
% =============================================================================
\chapter[Conclusions]{Conclusions and Recommendations}
\label{chap:conclusions}

\chapterquote{A thesis ends where a research programme begins.}{tradition
of the doctoral oral examination}

\begin{chapterpreview}
This chapter restates the principal contributions of the thesis,
summarises the conclusions that the bulk and per-particle evidence
jointly support, sets out the limitations of the present work, and
recommends three immediate avenues for follow-up: an instrumented
cyclic-loading extension; an in-track field-validation campaign on a
heavy-haul stretch; and a discrete-element model calibrated against
the paired-particle dataset reported in \cref{chap:morph-results}.
\end{chapterpreview}

\dropcap{T}{his} thesis set out to address a combined apparatus-and-characterisation
gap in the current assessment of railway ballast degradation under
drop-weight impact loading, specifically in the context of the \gls{RDSO}
\SIrange{20}{65}{\milli\metre} gradation envelope and the
\SI{22.5}{\tonne} Indian Railways axle-load stress range. The programme
comprised: the design, fabrication, and commissioning of a purpose-built
large-scale drop-weight impact apparatus at Gati Shakti Vishwavidyalaya;
the execution of a nine-configuration factorial test matrix crossing
three pad conditions (\gls{UR}~/~\gls{RUBM}~/~\gls{PUBM}) with two base
conditions (\gls{FB}~/~\gls{RB}) and two impact energies
(\SI{300}{\milli\metre}~/~\SI{450}{\milli\metre} drop); the
characterisation of bulk degradation through the Indian Railway Fouling
Index \FIIN{} on the post-test gradations; the paired three-dimensional
morphological characterisation of 45 individually-tracked tracer particles
through twenty-two descriptors before and after loading; and the
statistical analysis of both characterisations individually and in
synthesis. The present chapter states the principal conclusions of the
programme, enumerates the specific novel contributions, identifies the
limitations that bound the generalisability of the findings, and recommends
concrete directions for follow-up research.

\section{Conclusions}
\label{sec:concl-main}

Six principal conclusions are drawn from the experimental programme and
its analysis. They are stated here in the order in which they emerged
during the research.

\begin{finding}[title=Conclusion 1 -- Apparatus]
A large-scale drop-weight setup was successfully developed at \gls{GSV},
delivering consistent and reproducible impact energies (\SI{147}{\joule}
and \SI{221}{\joule}, corresponding to nominal contact stresses of
\SI{275}{\kilo\pascal} and \SI{523}{\kilo\pascal}) representative of
Indian Railways conditions.
\end{finding}

\begin{finding}[title=Conclusion 2 -- Settlement]
Ballast shows rapid initial settlement that stabilises after about 15
blows, confirming a 20-blow test as appropriate. Higher energy and rigid
bases cause greater settlement.
\end{finding}

\begin{finding}[title=Conclusion 3 -- Fouling]
Fouling varies significantly with test conditions~--- highest under rigid
base, no pad, high energy (\testlabel{T7}: \FIIN{} = \SI{3.386}{\percent});
lowest with polyurethane pad, flexible base, low energy (\testlabel{T1}:
\FIIN{} = \SI{0.533}{\percent}). A 6.4-fold span across the nine
configurations establishes that pad, base, and energy variables each
exert a material effect.
\end{finding}

\begin{finding}[title=Conclusion 4 -- Settlement vs.\ fouling]
Settlement and fouling measure different mechanisms~--- rearrangement vs.\
breakage~--- and are not interchangeable. \testlabel{T5} has the largest
settlement (\SI{48.6}{\milli\metre}) but only intermediate fouling;
\testlabel{T7} has the highest fouling but not the largest settlement.
Both metrics are required for full bulk-level assessment.
\end{finding}

\begin{finding}[title=Conclusion 5 -- Two damage modes]
\textbf{Mode~A (abrasion):} all 45 tracer particles undergo gradual
corner abrasion, increasing sphericity ($\Delta\Psi = +0.0041$ for
inliers, $p < 10^{-6}$) and reducing size across all
configurations.\par
\textbf{Mode~B (fracture):} catastrophic particle breakage occurs in
6/45 tracer particles, and \emph{all six are drawn from rigid-base
configurations}. The compliant sand base attenuates the impulsive force
enough to keep every tracer in the abrasion-only regime.
\end{finding}

\begin{finding}[title=Conclusion 6 -- Complementary characterisation]
\FIIN{} captures bulk fouling effects, while paired per-particle
morphology reveals damage mechanisms. Both are essential for a complete
characterisation: \FIIN{} alone cannot distinguish diffuse abrasion from
catastrophic fracture, and the per-particle protocol alone cannot speak
to assembly-level fines generation.
\end{finding}

\section{Novel contributions}
\label{sec:concl-novel}

\subsection{Apparatus and method}
\label{sec:concl-novel-apparatus}

\begin{enumerate}[label=(\arabic*),leftmargin=*,itemsep=4pt]
    \item A modified drop-weight impact apparatus for Indian ballast. A
          utility patent has been filed with the \gls{IPO}
          (\cref{chap:appB}).
    \item A standardised protocol calibrated to the Indian Railways
          stress window (\SIrange{275}{523}{\kilo\pascal}), recording
          per-blow peak force and five-milestone settlement---a full
          force-and-deformation history, versus the single post-test
          fines percentage of conventional Aggregate Impact Value testing.
    \item A morphological-analysis method (\cref{chap:appB}) that ingests
          paired pre/post \gls{STL} meshes and exports twenty-two
          descriptors per particle to a multi-sheet Excel workbook for
          direct comparison.
\end{enumerate}

\subsection{Scientific findings}
\label{sec:concl-novel-science}

\begin{enumerate}[label=(\arabic*),leftmargin=*,itemsep=4pt]
    \item Empirical demonstration that base condition (sand cushion vs.\
          rigid concrete) governs the selection between \gls{ModeA}
          (diffuse abrasion) and \gls{ModeB} (catastrophic fracture)
          under drop-weight impact loading of ballast, using a paired
          per-particle 3D morphological method. All six observed
          \gls{ModeB} events occurred on the rigid base.
    \item A heat-map of mean per-particle \%-change across fourteen
          descriptors for nine configurations: sign-uniformity in size
          and primary-shape descriptors, heterogeneity in shape-ratio
          descriptors, and sphericity and convexity as more robust
          diagnostics than angularity and roundness.
    \item Empirical demonstration that pad performance is energy- and
          base-dependent: PUBM dominates at the FB-300 and RB-450
          corners; RUBM edges ahead on RB-300. A single optimal pad
          cannot be specified independent of the substructure and
          loading energy.
    \item The first published nine-configuration factorial dataset
          pairing large-diameter drop-weight impact loading (\gls{RDSO}
          gradation, Indian Railways stress levels) with per-particle
          3D morphology---a reference for validating discrete-element
          models and machine-learning degradation predictors.
\end{enumerate}

\section{Limitations of the study}
\label{sec:concl-limitations}

Four limitations bound the generalisability of the conclusions and should
be recorded in any reading of the findings.

\begin{caveat}[title=Limitation 1 -- Sample size per configuration]
Five tracer particles supports the pooled \gls{ModeA} significance test
and \gls{ModeB} detection, but not per-configuration ANOVA or tight
\gls{ModeB} probability estimates. Expanding to \numrange{10}{20}
tracers is the single most effective follow-up improvement.
\end{caveat}

\begin{caveat}[title=Limitation 2 -- Single parent rock type]
Only basalt from one Vadodara quarry was tested. Parent rock affects
both \gls{ModeA} abrasion rate and \gls{ModeB} fracture threshold, so
generalisation requires extension to another \gls{RDSO}-approved rock.
\end{caveat}

\begin{caveat}[title=Limitation 3 -- Surface morphology only]
The EinScan structured-light scanner cannot detect sub-surface
micro-cracks, incipient fracture planes, or mineralogical alteration.
\gls{ModeA} particles may carry latent damage that accelerates
\gls{ModeB} fracture; X-ray CT would be required to detect this.
\end{caveat}

\begin{caveat}[title=Limitation 4 -- Laboratory scale]
The study is limited to a single mould diameter and a single hammer
mass. Field-scale validation, including the effect of variable train
speeds and long-term cyclic loading, is outside the present scope.
\end{caveat}

\section{Recommendations for future work}
\label{sec:concl-future}

Five concrete directions for follow-up research are identified, in rough
order of increasing scope. Each would usefully build on the methodological
platform established by the present thesis.

\begin{enumerate}[label=(\arabic*),leftmargin=*,itemsep=4pt]
    \item \textbf{Expanded tracer-count study.} Repeat the present
          nine-configuration programme with \numrange{15}{20} tracer
          particles per configuration to support per-configuration
          ANOVA and tighter estimates of \gls{ModeB} event probabilities.
    \item \textbf{Multi-rock-type extension.} Extend the programme to
          granite and quartzite ballast, using the same nine-configuration
          matrix and the same per-particle morphology pipeline, to map
          the parent-rock dependence of the \gls{ModeA}~/~\gls{ModeB}
          crossover.
    \item \textbf{AI/ML-based prediction of degradation behaviour.} Use
          the per-particle paired morphological dataset to train
          supervised machine-learning models. A trained surrogate would
          allow rapid screening of pad--base combinations without
          running the full impact campaign, and feature-importance
          analysis would identify which morphological descriptors most
          strongly govern degradation outcomes.
    \item \textbf{Cyclic-plus-impact combined loading.} Develop a
          composite loading method that applies cyclic-triaxial loading
          (mimicking ordinary traffic) interrupted periodically by
          drop-weight impact events (mimicking transition-zone hot-spots),
          and report the corresponding
          \gls{ModeA}~/~\gls{ModeB} evolution. This is the closest
          laboratory analogue to real service loading.
    \item \textbf{Field correlation.} Couple the laboratory degradation
          metrics with in-situ measurements at instrumented bridge
          approaches, rail joints, or transition zones on Indian
          Railways, to calibrate the laboratory predictions against
          observed track-geometry degradation rates.
\end{enumerate}
